\chapter{Team Responsibilities and Contributions}

The assignment explicitly evaluates not just technical output, but also whether the team had clear
ownership boundaries and a believable plan for cross-functional work. The project was therefore
organized around service ownership with shared decision-making for cross-cutting flows such as AI,
permissions, realtime collaboration, and versioning.

\section{Ownership Model}

\begin{table}[H]
\centering
\caption{Primary ownership and contribution areas}
\label{tab:ownership}
\begin{tabularx}{\textwidth}{p{2.4cm}p{4.4cm}X}
\toprule
\textbf{Member} & \textbf{Primary Ownership} & \textbf{Concrete Responsibilities} \\
\midrule
Tabina & Architecture and infrastructure & Architectural drivers, C4 level 1 and level 2 diagrams, feature decomposition, repository organization, methodology, milestone planning, final repository consistency \\
Noel & Backend and API layer & Document management requirements, user management requirements, latency/scalability targets, C4 level 3 backend component view, API contracts, auth/authorization design, realtime feasibility, backend MVP implementation \\
Nurtore & AI and data model & AI requirements, quota/failure handling, privacy and security NFRs, AI integration deep design, AI model strategy ADR, AI/data-model relationships \\
Ming Ming & Frontend and UX & Collaboration requirements, usability NFRs, user stories, frontend PoC/MVP implementation, API usability validation, role-aware UX behavior \\
\bottomrule
\end{tabularx}
\end{table}

\section{Cross-Cutting Collaboration}

Several major features deliberately crossed ownership boundaries.

\begin{itemize}
  \item \textbf{Realtime collaboration} required coordination between Ming Ming's frontend work,
  Noel's backend contracts and permission model, and the dedicated realtime service design.
  \item \textbf{AI assistance} touched frontend UX, backend policy enforcement, the AI provider
  abstraction, the data model, and versioning logic.
  \item \textbf{Permissions and sharing} required agreement between the frontend role UX, backend
  RBAC enforcement, and realtime write restrictions.
  \item \textbf{Final demo reliability} required infrastructure-oriented work such as shared runtime
  configuration, root scripts, Docker Compose, CI, and deterministic reset behavior.
\end{itemize}

To avoid fragmented decision making, the team treated backend contracts, realtime behavior, and AI
interaction flows as interface boundaries that had to be reviewed jointly. This mattered because the
assignment's PoC requirement was not just “make something run,” but “make the code reflect the
architecture document.” Shared ownership therefore meant aligning contracts and assumptions, not
blurring responsibility.

\section{Decision-Making Process}

Technical disagreements were resolved with a bias toward the smallest correct design that preserved
clarity, demo reliability, and grading alignment. In practice, the decision rule was:

\begin{enumerate}
  \item identify the architectural driver being optimized (correctness, latency, maintainability,
  cost, or demo reliability),
  \item prefer the approach that reduced implementation risk while still satisfying the assignment,
  \item record the choice through an ADR when the trade-off affected multiple parts of the system.
\end{enumerate}

This is why the delivered MVP uses local bearer-token authentication rather than a full OAuth/OIDC
 integration, SQLite rather than a larger operational database deployment, and a provider abstraction
with a deterministic stub AI path for demos. These choices intentionally optimize feasibility and
traceability without abandoning the target architecture described in the report.
